% ===================================================================
%  పట్టిక సంఖ్య 6 — ఇంటి లోపలి భాగం, సామానుసుమాను, ఆహారం, వెలుతురు.
%
%  Traduction, faite en 2026, en telougou d'aujourd'hui, sur l'ido de
%  texto/io. Regles de la colonne : voir 10-pattika-01.tex. Cinq
%  scenes, cinq numerotations : les cles portent c1 a c5.
%
%  DEUX NOTES SUR LA MEME PAGE, DEUX \VUnotes. Le fac-simile les
%  compose dans un seul bloc separe par \par, mais elles portent
%  deux cles : les ecrire ensemble ferait disparaitre la seconde.
%
%  ECRIT AVEC LA LISTE DE parigi.py EN MAIN, comme le tableau 5.
%  Quarante-six paires signalees, dont une vingtaine sont de vrais
%  rapports modificateur-tete : la corbeille A OUVRAGE PRES DU
%  TABOURET, l'oreiller QUI EST AVEC L'EDREDON, le grand cadre
%  PENDU AU MUR, le seau A CHARBON PLEIN DE HOUILLE QUE LA CUISINIERE
%  A REMPLI, le baquet OU L'ON LAVE LES LEGUMES, le tonneau MUNI DE
%  SON ROBINET. Le reste de la liste — « ornee d'une statuette, de
%  deux vases » — sont des coordinations, que le telougou rend dans
%  l'ordre : l'outil sur-signale, et c'est le bon sens du biais.
%
%  LE TELOUGOU REPOND AUX DEUX NOTES, comme le marathi y avait
%  repondu. La premiere explique que « babo » vient de l'anglais
%  baby, du francais bebe, de l'italien bambino ; le telougou dit
%  పసిబిడ్డ. La seconde aligne cinq langues pour le « lambrequino »,
%  cette bande d'etoffe qui pend au-dessus d'une porte ou d'une
%  cheminee ; le telougou dit తోరణం — le mot de la guirlande qu'on
%  suspend au linteau les jours de fete, et qui est exactement
%  l'objet. Quatre notes du livret, quatre mots que ni le marathi ni
%  le telougou n'ont eu a emprunter.
%
%  LA CUISINE EST LE MEILLEUR MORCEAU, COMME EN MARATHE, et pour la
%  meme raison : c'est la piece ou une langue a le plus servi.
%
%    పొయ్యి (1)    le foyer        నిప్పు (3)    le feu
%    కొలిమితిత్తి (8) le soufflet    లాంతరు (11)  la lanterne
%    గరిటె (21)    la louche       మూకుడు (47)   la poele
%    కుండలు (43)   les pots        జల్లెడ (28)   la passoire
%    తురుము పీట (29) la rape        గరాటు (38)   l'entonnoir
%    తొట్టి (49)   l'evier         బిందె (50)    la grande cruche
%    సొరుగు (33)   le tiroir       చీపురు (34)   le balai
%    కూర (24)      la sauce        అర (42)      la planche
%    అగ్గిపెట్టె (13) la boite d'allumettes
%
%  ET « కూర » EST EXACTEMENT LA SAUCE DE L'IDO, ce qui n'est pas un
%  hasard : le mot telougou designe le plat mouille qui accompagne le
%  riz, c'est-a-dire la meme chose que le livret appelle sauco et que
%  la cuisiniere remue pour qu'elle ne brule pas.
%
%  LA CHAMBRE DE MEME : మంచం le lit, తలగడ l'oreiller, పరుపు le
%  matelas, దుప్పటి le drap, వితానం le ciel de lit, బీరువా l'armoire,
%  కూజా le broc, దువ్వెన le peigne, కత్తెర les ciseaux, సూదులు les
%  aiguilles, దారం le fil, తాళంచెవి la clef, అద్దం le miroir, ఊయల le
%  berceau.
% ===================================================================

%%K t06-apar-1 apar
\VUsaut{6.0mm}
\VUtitre{15.0pt}{60}{పట్టిక సంఖ్య 6}
\VUsaut{1.4mm}
\VUtitre{11.4pt}{0}{ఇంటి లోపలి భాగం, సామానుసుమాను, ఆహారం,}
\VUsaut{0.4mm}
\VUtitre{11.4pt}{0}{వెలుతురు.}
\VUsaut{2.2mm}

%%K t06-apar-2 apar
ఇల్లు పూర్తయింది. జాన్ మామయ్య తన కొత్త ఇంట్లో కాపురం పెట్టారు, ఆ ఇంటి అమరిక మనకు తెలుసు.
ఇప్పుడు ఆయన దాన్ని ఎలా సర్దుకున్నారో చూద్దాం. ముందుగా ఇది చూద్దాం :

%%K t06-tit-1 sub
\VUpk{10.2pt}{40}{పడక గది.}
\VUsaut{3.0mm}

%%K t06-c1-01-1 p
1. --- మనం సరైన సమయంలో రాలేదు, ఎందుకంటే పనిమనిషి \VUgras{గది} సర్దడంలో మునిగి ఉంది.

%%K t06-c1-02-1 p
2. --- \VUgras{ప్రసాధన బల్లపై} \textsuperscript{(1)} అక్కడక్కడా రకరకాల వస్తువులు ఉన్నాయి,
వాటితో మనిషి కడుక్కుంటాడు, తల దువ్వుకుంటాడు — మొత్తానికి తయారవుతాడు. అవి \VUgras{గంగాళం}
\textsuperscript{(2)} మరియు \VUgras{నీళ్ళ కూజా} \textsuperscript{(3)}, \VUgras{జుట్టు బ్రష్}
\textsuperscript{(4)}, \VUgras{దువ్వెన} \textsuperscript{(5)}, \VUgras{సబ్బు}
\textsuperscript{(6)} మరియు \VUgras{స్పాంజి} \textsuperscript{(7)}, చివరగా ద్రవ దంతచూర్ణపు
\VUgras{చిన్న సీసా} \textsuperscript{(8)}.

%%K t06-c1-03-1 p
3. --- నేలపై, ప్రసాధన బల్ల ముందు, ఇదిగో మురికినీళ్ళు పోసుకునే \VUgras{బకెట్}
\textsuperscript{(9)}.

%%K t06-c1-04-1 p
4. --- \VUgras{బయటి గదిలో} \textsuperscript{(10)} మనకు కొన్ని \VUgras{గుడ్డల కొక్కేలు}
\textsuperscript{(13)} మరియు రెండు \VUgras{గొడుగులు} \textsuperscript{(11)} కనిపిస్తాయి —
\VUgras{గొడుగుల నిలువులో} \textsuperscript{(12)}.

%%K t06-c1-05-1 p
5. --- తలుపుకు అవతలివైపు \VUgras{అద్దాల బీరువా} \textsuperscript{(14)} ఉంది, దానిలో
\VUgras{బట్టలు} \textsuperscript{(15)} ఉంటాయి.

%%K t06-c1-06-1 p
6. --- \VUgras{పనిమనిషి} \textsuperscript{(6)} \VUgras{మంచం} \textsuperscript{(17)}
సర్దుతున్నప్పుడే దాని వేర్వేరు భాగాలు చూసేద్దాం. \VUgras{మంచపు చట్రంలో} \textsuperscript{(20)}
మంచి \VUgras{స్ప్రింగుల పరుపు} \textsuperscript{(21)} ఉంది, దానిపై యుస్తీనా వెంటనే
\VUgras{ఉన్ని పరుపు} \textsuperscript{(22)} వేస్తుంది. ఆ తర్వాత కుర్చీల మీది రెండు
\VUgras{దుప్పట్లు} \textsuperscript{(23)} మరియు \VUgras{కప్పుకునే దుప్పటి}
\textsuperscript{(24)} తీసుకుంటుంది. అప్పుడు మంచం తలవైపు \VUgras{దీర్ఘ దిండు}
\textsuperscript{(25)} మరియు \VUgras{తలగడ} \textsuperscript{(26)} పెడుతుంది — అవి
\VUgras{మెత్తని దుప్పటితో} \textsuperscript{(27)} కలిసి \VUgras{పైదుప్పటిపై}
\textsuperscript{(32)} ఉన్నాయి. \VUgras{తెరలు} \textsuperscript{(19)} మరియు
\VUgras{మంచపు వితానం} \textsuperscript{(18)} దులిపేందుకు ఆమె ఈకల చీపురు వాడుతుంది, ఇదిగో పనిలో
కొంత భాగం అయిపోయింది.

%%K t06-c1-07-1 p
7. --- ఆ తర్వాత ఆమె \VUgras{మంచం పక్క బల్ల} \textsuperscript{(28)} కొంచెం సర్దాలి,
\VUgras{అలారం గడియారానికి} \textsuperscript{(29)} కీ ఇవ్వాలి, \VUgras{రాత్రి దీపం}
\textsuperscript{(30)} సిద్ధం చేయాలి, స్టాండు శుభ్రం చేసేందుకు \VUgras{కొవ్వొత్తి}
\textsuperscript{(31)} తీయాలి.

%%K t06-tit-2 sub
\VUsaut{5.0mm}
\VUpk{10.2pt}{40}{కుట్టుపని బుట్ట మరియు పొయ్యి సామగ్రి.}
\VUsaut{3.0mm}

%%K t06-c1-08-1 p
8. --- \VUgras{కుట్టుపని బుట్టను} \textsuperscript{(34)} కూడా సర్దాల్సిందే — అది \VUgras{పీట}
\textsuperscript{(33)} దగ్గర ఉంది. ఆమె \VUgras{అల్లిక పని} \textsuperscript{(35)} మడిచి
పెట్టాలి, \VUgras{కుట్టుపని బల్లలో} \textsuperscript{(38)} \VUgras{వేలికప్పు}, \VUgras{దారం}
\textsuperscript{(36)}, \VUgras{సూదులు} \textsuperscript{(37)} మరియు \VUgras{కత్తెర}
\textsuperscript{(39)} పెట్టాలి, \VUgras{తాళంచెవితో} \textsuperscript{(40)} బల్ల మూత వేయాలి.

%%K t06-c1-09-1 p
9. --- మధ్యలోని \VUgras{ఒంటికాలి బల్లపై} \textsuperscript{(41)} యుస్తీనా
\VUgras{గాజు కూజాలోని} \textsuperscript{(42)} నీళ్ళు మారుస్తుంది. \VUgras{పంచదార డబ్బాలో}
\textsuperscript{(43)} \VUgras{పంచదార} వేస్తుంది, \VUgras{పంచదార పట్టకారు}
\textsuperscript{(44)} తుడుస్తుంది, \VUgras{బట్టల బ్రష్} \textsuperscript{(46)} తీసి
పెడుతుంది.

%%K t06-c1-10-1 p
10. --- గది కుడివైపు ఆమెకు తక్కువ పని పెడుతుంది, ఎందుకంటే \VUgras{పొడవాటి ఆరామకుర్చీ}
\textsuperscript{(47)} మరియు దాని \VUgras{దిండు} \textsuperscript{(48)} సరైన చోటే ఉన్నాయి, చలి
లేకపోవడంతో \VUgras{పొయ్యి} \textsuperscript{(49)} సామగ్రిలో ఏదీ కదల్లేదు. \VUgras{పార}
\textsuperscript{(50)}, \VUgras{పట్టకారు} \textsuperscript{(51)},
\VUgras{కట్టెల ఇనుప గుర్రాలు} \textsuperscript{(55)}, \VUgras{బూడిద అడ్డుకట్ట}
\textsuperscript{(56)} శుభ్రంగా ఉన్నాయి; \VUgras{మంట తెర} \textsuperscript{(54)} కదల్లేదు,
\VUgras{కట్టెల పెట్టె} \textsuperscript{(52)} బాగా నిండి ఉంది — \VUgras{కట్టెలతో}
\textsuperscript{(53)}.

%%K t06-noto-1 noto
\VUnotes{143.2mm}{%
(*) Babo = పసిబిడ్డ; ఇంగ్లిష్ \textit{baby}, ఫ్రెంచ్ \textit{bébé}, ఇటాలియన్
\textit{bambino}.%
}

%%K t06-noto-2 noto
\VUnotes{143.2mm}{%
(*) Lambrequino = ఫ్రెంచ్ \textit{lambrequin}, స్పానిష్ \textit{lambrequines}, పోర్చుగీసు
\textit{lambrequins}, జర్మన్ మరియు ఇంగ్లిష్‌లోనూ \textit{lambrequins}; అంటే జ. ఇం. ఫ్రెం. పో.
స్పా.%
}

%%K t06-c1-11-1 p
11. --- అయినా ఆమె \VUgras{లోలకపు గడియారం} \textsuperscript{(57)}, \VUgras{దీపస్తంభాలు}
\textsuperscript{(58)} మరియు \VUgras{జేబు సామాను పళ్ళేలు} \textsuperscript{(59)} దులపాలి,
చివరగా \VUgras{అద్దం} \textsuperscript{(60)} తుడవాలి.

%%K t06-c1-12-1 p
12. --- పసివాడి (\VUgras{పసిబిడ్డ} (*)) \VUgras{ఊయల} \textsuperscript{(61)} విషయానికొస్తే, అది
ఇప్పటికే సిద్ధంగా ఉంది.

%%K t06-tit-3 sub
\VUsaut{5.0mm}
\VUpk{10.2pt}{40}{సావిడి.}
\VUsaut{3.0mm}

%%K t06-c2-01-1 p
1. --- ఇప్పుడు ఇదిగో సావిడి. ఇది చాలా అందమైన \VUgras{గది}. కుడి మూలలో ఉన్న పొయ్యిని ఒక
సున్నితమైన \VUgras{విగ్రహం} \textsuperscript{(1)}, రెండు అందమైన \VUgras{పూలదాన్లు}
\textsuperscript{(3)} అలంకరిస్తున్నాయి, దానిపై మఖమలు \VUgras{తోరణం} \textsuperscript{(2)} (*)
వేలాడుతోంది.

%%K t06-c2-02-1 p
2. --- పొయ్యిలోని నిప్పురవ్వలు తివాచీని కాల్చకుండా ఉండేందుకు పొయ్యి ముందు రాగి
\VUgras{నిప్పురవ్వల జాలి} \textsuperscript{(4)} పెట్టారు.

%%K t06-c2-03-1 p
3. --- దగ్గరలో ఒక చక్కని ఆడవాళ్ళ \VUgras{రాత బల్ల} \textsuperscript{(5)} తెరిచి ఉంది. దానిపై
\VUgras{ఇల్లాలు} \textsuperscript{(23)} తన ఉత్తరాలు రాస్తుంది.

%%K t06-c2-04-1 p
4. --- కిటికీకి \VUgras{కిటికీ తెరలు} \textsuperscript{(6)} ఉన్నాయి, వాటికి
\VUgras{కట్టుతాళ్ళు} \textsuperscript{(8)} ఉన్నాయి, ఆ తాళ్ళు \VUgras{కొక్కేలకు}
\textsuperscript{(7)} తగిలించి ఉన్నాయి; \VUgras{పైనుంచి వేలాడే} \textsuperscript{(9)} మందపాటి
గుడ్డ తెరలు కూడా ఉన్నాయి.

%%K t06-c2-05-1 p
5. --- కిటికీ గూట్లో ఒక \VUgras{కుండీలో} \textsuperscript{(11)} పచ్చని \VUgras{మొక్క}
\textsuperscript{(10)} ఉంది — ఒక అందమైన తాటిమొక్క, దాని ఆకులు దాదాపు \VUgras{కిరసనాయిలు దీపం}
\textsuperscript{(12)} వరకు చాచుకున్నాయి.

%%K t06-c2-06-1 p
6. --- దీపపు \VUgras{గొడుగును} \textsuperscript{(13)} మారియా సర్దింది — ఆ ముద్దుల
\VUgras{యువతి} \textsuperscript{(14)}, ఆమె \VUgras{పియానో} \textsuperscript{(15)} దగ్గర ఉంది.
చాలా నిటారుగా \VUgras{పీటపై} \textsuperscript{(16)} కూర్చుని ఆమె సంగీతం అభ్యసిస్తోంది. ఆమె
చాలా నేర్పరి, నీటుగా ఉంటుంది — అది ఆమె \VUgras{సంగీత బీరువా} \textsuperscript{(17)} చూస్తే
తెలుస్తుంది, అది చక్కగా సర్ది ఉంది.

%%K t06-tit-4 sub
\VUsaut{5.0mm}
\VUpk{10.2pt}{40}{కొంటె కుర్రాడు.}
\VUsaut{3.0mm}

%%K t06-c2-07-1 p
7. --- ఆమె తమ్ముడు పౌలు ఒక \VUgras{పుస్తకం} \textsuperscript{(19)} వెతుకుతున్నాడు,
\VUgras{పుస్తకాల బీరువాలో} \textsuperscript{(18)}. అతను \VUgras{యువకుడు}
\textsuperscript{(20)}, చంచలుడు, భరించలేనివాడు. చాలా పైనున్న పుస్తకం అందుకునేందుకు బీరువాకు
వేలాడుతూ, పైన ఉన్న \VUgras{అర్ధప్రతిమను} \textsuperscript{(21)} పడేసే ప్రమాదం అతను
తెచ్చుకుంటున్నాడు.

%%K t06-c2-08-1 p
8. --- ఆ కొంటెతనం చేసేందుకు అతనికి ఈ అవకాశం దొరికింది, ఎందుకంటే అతని తల్లి ఒక స్నేహితురాలితో —
కొన్ని రోజులు ఉండేందుకు వచ్చిన ఒక \VUgras{మహిళతో} \textsuperscript{(22)} — మాట్లాడటంలో మునిగి
ఉంది. తల్లి \VUgras{సోఫాపై} \textsuperscript{(24)} కూర్చుంది, \VUgras{ఆరామకుర్చీ}
\textsuperscript{(25)} దగ్గర; ఆమె స్నేహితురాలు ఆ \VUgras{కుర్చీపై} \textsuperscript{(27)}
ఉంది, దాని \VUgras{వీపు ఆనుపు} \textsuperscript{(26)} మాత్రమే కనిపిస్తోంది.

%%K t06-c2-09-1 p
9. --- పెద్ద \VUgras{చట్రంలో} \textsuperscript{(30)} — అది గోడకు వేలాడుతోంది — ఒక బంధువు
స్త్రీ \VUgras{చిత్రపటం} \textsuperscript{(29)} ఉంది. బల్లపై పెట్టిన \VUgras{ఆల్బంలో}
\textsuperscript{(32)} మిగతా ఇంటివాళ్ళ \VUgras{ఛాయాచిత్రాలు} \textsuperscript{(33)} పోగుచేసి
ఉన్నాయి. పిల్లలు దాన్ని ఇప్పుడే తిరగేశారు, ఎందుకంటే \VUgras{కుర్చీలు} \textsuperscript{(31)}
ఇంకా బల్ల చుట్టూ చేరి ఉన్నాయి.

%%K t06-c2-10-1 p
10. --- పింగాణీ \VUgras{చైనా పూలదాని} \textsuperscript{(34)} — అది \VUgras{కాగితం కోసే కత్తి}
\textsuperscript{(35)} దగ్గర ఉంది — మారియాకు ఆమె పండుగనాడు కానుకగా వచ్చింది; గది మూలను
అలంకరించే \VUgras{మడత తెర} \textsuperscript{(36)} ఆమె మామయ్య చైనా నుంచి తెచ్చారు.

%%K t06-c2-11-1 p
11. --- అందమైన \VUgras{చిత్రాలతో} \textsuperscript{(37)} — అవి \VUgras{గోడకు}
\textsuperscript{(42)} తగిలించి ఉన్నాయి — ఆనందంగా, \VUgras{వేలాడే గుత్తిదీపంతో}
\textsuperscript{(38)} వెలుగుతూ — దాని \VUgras{విద్యుత్ దీపాలు} \textsuperscript{(39)}
సాయంత్రం ఆనందంగా మెరుస్తాయి — ఆ సావిడి చాలా హాయిగా కనిపిస్తుంది. చెక్క నేలపై పరిచిన మెత్తని
\VUgras{తివాచీ} \textsuperscript{(40)} మరియు అంత సౌకర్యవంతమైన \VUgras{విద్యుత్ గంట}
\textsuperscript{(41)} దాని సౌఖ్యాన్ని పూర్తి చేస్తాయి.

%%K t06-tit-5 sub
\VUsaut{3.6mm}
\VUpk{10.2pt}{40}{భోజన గది.}
\VUsaut{3.0mm}

%%K t06-c3-01-1 p
1. --- భోజనపు గంట మోగింది. మారియా తప్ప ఇల్లంతా బల్ల చుట్టూ చేరింది. చక్కని పింగాణీ
\VUgras{వేడిపొయ్యి} \textsuperscript{(1)} భోజన గదిని హాయిగా వెచ్చబెట్టింది — దానికి సన్నని
\VUgras{గొట్టం} \textsuperscript{(2)} ఉంది.

%%K t06-c3-02-1 p
2. --- ఈ గదిలో సామాను ఎక్కువ లేదు. గోడ వెంట, \VUgras{చిన్న బల్లపైన} \textsuperscript{(4)}, ఒక
అలంకారపు \VUgras{నీటిబుగ్గ} \textsuperscript{(3)} వేలాడుతోంది. చాలా దగ్గరలో
\VUgras{సామాన్ల బీరువా} \textsuperscript{(5)} ఉంది, దానిపై చల్లని పదార్థాలు, తీపి వంటకాల
పళ్ళేలు, వెంటనే అవసరం లేని బల్ల సామాను పెడతారు.

%%K t06-c3-03-1 p
3. --- ఆ విధంగా పై అరపై మనకు \VUgras{కాల్చిన కోడిపిల్ల} \textsuperscript{(7)}, \VUgras{చేప}
\textsuperscript{(8)} మరియు ఒక \VUgras{పళ్ళెం} \textsuperscript{(10)} కనిపిస్తాయి — అది
\VUgras{పండ్లతో} \textsuperscript{(9)} నిండి ఉంది. కింద ఇదిగో \VUgras{చీజ్}
\textsuperscript{(11)}, \VUgras{రొట్టె} \textsuperscript{(12)}, మరియు \VUgras{సీసాల పళ్ళెం}
\textsuperscript{(13)} — దానిలో కోసంబరికి కావలసినవన్నీ ఉంటాయి : నూనె, వెనిగర్, \VUgras{ఉప్పు}
\textsuperscript{(14)} మరియు \VUgras{మిరియాలు} \textsuperscript{(15)}. చివరగా, కింది అరపై
\VUgras{కేకు} \textsuperscript{(17)} ఉంది, \VUgras{ఆవాల డబ్బా} \textsuperscript{(16)} దగ్గర.

%%K t06-c3-04-1 p
4. --- భోజన గదిలోని గడియారాన్ని \VUgras{గోడ గడియారం} \textsuperscript{(20)} అంటారు, దాని ఆకారం
చాలా ప్రత్యేకం. అది ఒక చెక్క చట్రంలో అమర్చి ఉంది, ఆ చట్రంలో ఇంకా \VUgras{వాయుభారమాపకం}
\textsuperscript{(19)} మరియు \VUgras{ఉష్ణమాపకం} \textsuperscript{(18)} ఉన్నాయి.

%%K t06-tit-6 sub
\VUsaut{4.6mm}
\VUpk{10.2pt}{40}{భోజన వడ్డన.}
\VUsaut{3.0mm}

%%K t06-c3-05-1 p
5. --- గది గోడలన్నిటికీ మైనం పూసిన ఓక్ \VUgras{చెక్క పలకలు} \textsuperscript{(21)} బిగించి
ఉన్నాయి. వరండాలోకి వెళ్ళే తలుపును గుడ్డ \VUgras{తలుపు తెర} \textsuperscript{(22)} బాగానే
మూస్తుంది. భోజనం తర్వాత ఇల్లంతా అక్కడ కాఫీ తాగేందుకు వెళ్తుంది. \VUgras{కప్పులు}
\textsuperscript{(24)} మరియు \VUgras{సాసర్లు} \textsuperscript{(25)} ఇప్పటికే సిద్ధంగా ఉన్నాయి
— \VUgras{వడ్డన బల్లపై} \textsuperscript{(23)}.

%%K t06-c3-06-1 p
6. --- బల్లపైన పైకప్పుకు \VUgras{వేలాడే దీపం} \textsuperscript{(26)} బిగించి ఉంది, దాని చాలా
ప్రకాశవంతమైన దీపం సాయంత్రం భోజన గది అంతటినీ వెలిగిస్తుంది.

%%K t06-c3-07-1 p
7. --- ఇప్పుడు \VUgras{భోజన బల్లనే} \textsuperscript{(27)} చూద్దాం. ఇల్లంతా లోపలికి
వచ్చినప్పుడు వడ్డన సిద్ధంగా ఉంది. \VUgras{బల్లగుడ్డపై} \textsuperscript{(28)} ప్రతి భోజనం
చేసేవాడి చోట \VUgras{పళ్ళెం} \textsuperscript{(33)}, \VUgras{చేతిగుడ్డ}
\textsuperscript{(29)}, \VUgras{చెంచా} \textsuperscript{(34)}, \VUgras{ఫోర్కు}
\textsuperscript{(35)}, \VUgras{కత్తి} \textsuperscript{(36)} మరియు \VUgras{గ్లాసు}
\textsuperscript{(32)} పెట్టి ఉన్నాయి. ద్రాక్షసారాకు \VUgras{సీసాలు} \textsuperscript{(30)},
నీళ్ళకు \VUgras{కూజా} \textsuperscript{(31)} కూడా ఉన్నాయి.

%%K t06-c3-08-1 p
8. --- \VUgras{నౌకరు} \textsuperscript{(39)} ముందుగా \VUgras{సూపు గిన్నె}
\textsuperscript{(37)} తెచ్చాడు, దానిలో ఇంకా కొంత సూపు ఆవిరులు కక్కుతోంది. ఇప్పుడు అతను మొదటి
\VUgras{వంటకం} \textsuperscript{(38)} వడ్డిస్తున్నాడు — మాంసపు వంటకం. ఆ తర్వాత మనం ఇప్పుడే
పేర్కొన్నవి అందిస్తాడు, చివరగా, \VUgras{తీపి వంటకం} తర్వాత, యజమానులు వరండాలోకి వెళ్ళిన సందు
చూసుకుని బల్ల సామాను తీసేసి భోజన గది మళ్ళీ సర్దుతాడు.

%%K t06-c3-08-2 p
\textit{గమనిక.} ---
\textit{ఆహారానికి సంబంధించిన రోజువారీ పదాలు ఇక్కడ చాలా క్లుప్తంగా ఇచ్చాం. ఆ జాబితా « ఉపాహారశాల » (సంఖ్య 13) మరియు « సంత » (సంఖ్య 15) అనే రెండు పట్టికల్లో పూర్తి చేయబడుతుంది.}

%%K t06-tit-7 sub
\VUsaut{4.6mm}
\VUpk{10.2pt}{40}{వంటగది.}
\VUsaut{3.0mm}

%%K t06-c4-01-1 p
1. --- సగం తెరిచిన తలుపు గుండా మనం చూస్తున్న వంటగదిలో ఒక చిత్రమైన చిందరవందర కనిపిస్తుంది.
\VUgras{పొయ్యి} \textsuperscript{(1)} ముందు — దానిలో \VUgras{నిప్పు} \textsuperscript{(3)}
చిటపటలాడుతోంది — \VUgras{కబాబు కడ్డీ యంత్రం} \textsuperscript{(5)} అమర్చి ఉంది. చక్కని
\VUgras{కాల్చిన మాంసం} \textsuperscript{(7)} \VUgras{కడ్డీకి గుచ్చి} \textsuperscript{(6)}
ఉంది, అది కాలుతోంది.

%%K t06-c4-02-1 p
2. --- ఇప్పుడే వాడిన \VUgras{కొలిమితిత్తి} \textsuperscript{(8)} మరియు \VUgras{బొగ్గుల పార}
\textsuperscript{(9)} \VUgras{కట్టెలలాగే} \textsuperscript{(10)} నేలపైనే ఉండిపోయాయి.

%%K t06-c4-03-1 p
3. --- పొయ్యి పైభాగం సర్ది ఉంది; \VUgras{లాంతరు} \textsuperscript{(11)}, \VUgras{అగ్గిపెట్టె}
\textsuperscript{(13)} — దానిలో \VUgras{అగ్గిపుల్లలు} \textsuperscript{(12)} ఉన్నాయి —,
\VUgras{కొవ్వొత్తి స్టాండు} \textsuperscript{(14)} మరియు \VUgras{కొవ్వొత్తి పీట}
\textsuperscript{(15)} — దానిపై \VUgras{కొవ్వొత్తి} \textsuperscript{(16)} ఉంది — అన్నీ తమ చోట
ఉన్నాయి. కానీ పెద్ద \VUgras{బొగ్గుల బకెట్} \textsuperscript{(18)}, \VUgras{రాతిబొగ్గుతో}
\textsuperscript{(17)} నిండినది — దాన్ని \VUgras{వంటమనిషి} \textsuperscript{(23)} ఇప్పుడే
నేలమాళిగలో నింపుకొచ్చింది — వంటగది మధ్యలోనే ఉండిపోయింది.

%%K t06-c4-04-1 p
4. --- నిజానికి ఆ పాపం ఆవిడ హడావిడి పడాలి, అప్పుడే భోజనం సమయానికి సిద్ధమవుతుంది. ఇప్పుడు ఆమె
\VUgras{వంట పొయ్యి} \textsuperscript{(19)} దగ్గర ఉంది, ఎక్కువ మాడిపోకూడని \VUgras{కూర}
\textsuperscript{(24)} కలుపుతోంది.

%%K t06-c4-05-1 p
5. --- \VUgras{ఓవెన్‌లో} \textsuperscript{(20)} ఒక కేకు కాలుతోంది, పిల్లల చిరుతిండి కోసం ఆమె
దాన్ని తయారుచేసింది. వంట పొయ్యి పైన \VUgras{గరిటె} \textsuperscript{(21)} మరియు
\VUgras{జల్లెడ గరిటె} \textsuperscript{(22)} వేలాడుతున్నాయి.

%%K t06-tit-8 sub
\VUsaut{4.0mm}
\VUpk{10.2pt}{40}{వంట సామాను.}
\VUsaut{3.0mm}

%%K t06-c4-06-1 p
6. --- గది చివర తగిలించిన \VUgras{వంట సామాను} \textsuperscript{(25)} చాలా శుభ్రంగా ఉంది.
అందమైన \VUgras{రాగి గిన్నెలు} \textsuperscript{(26)}, \VUgras{పాల గిన్నె}
\textsuperscript{(27)}, \VUgras{జల్లెడ} \textsuperscript{(28)} మరియు \VUgras{తురుము పీట}
\textsuperscript{(29)} జాగ్రత్తగా తోమి ఉన్నాయి.

%%K t06-c4-07-1 p
7. --- \VUgras{ఉప్పు డబ్బా} \textsuperscript{(30)} సామాన్ల బీరువాపై \VUgras{టీ కూజా}
\textsuperscript{(31)} దగ్గర, \VUgras{నూనె సీసా} \textsuperscript{(32)} దగ్గర పెట్టి ఉంది.
\VUgras{సొరుగులో} \textsuperscript{(33)} బహుశా భోజన సామానూ వంటగది కత్తులూ ఉంటాయి.

%%K t06-c4-08-1 p
8. --- \VUgras{చీపురు} \textsuperscript{(34)} మరియు \VUgras{ఈకల చీపురు} \textsuperscript{(35)}
మూలలో ఉన్నాయి. వంటమనిషి ఇప్పుడే \VUgras{చెక్క దిమ్మె} \textsuperscript{(36)} వాడింది, దాన్ని
కిటికీ దగ్గర వదిలేసింది.

%%K t06-c4-09-1 p
9. --- \VUgras{చేతి తువ్వాలు} \textsuperscript{(37)} గోడకు వేలాడుతోంది, \VUgras{గరాటు}
\textsuperscript{(38)} దగ్గర. వంటగదిలోని ఆ భాగంలో \VUgras{కాల్చే జల్లి}
\textsuperscript{(39)}, \VUgras{నరికే కత్తి} \textsuperscript{(40)} మరియు \VUgras{నరికే పలక}
\textsuperscript{(41)} పెట్టి ఉన్నాయి. ఆ తర్వాత, నరికే కత్తి పైన, ఒక \VUgras{అరపై}
\textsuperscript{(42)} \VUgras{కుండలు} \textsuperscript{(43)}, \VUgras{మరిగించే గిన్నె}
\textsuperscript{(44)} — దాని \VUgras{మూతతో} \textsuperscript{(45)} — మరియు
\VUgras{పెద్ద గిన్నె} \textsuperscript{(46)} ఉన్నాయి. \VUgras{మూకుడు} \textsuperscript{(47)}
దగ్గరే వేలాడుతోంది.

%%K t06-c4-10-1 p
10. --- ఈ మూలలో రాగి \VUgras{కుళాయి} \textsuperscript{(48)} మాత్రమే మెరుస్తోంది.
\VUgras{తొట్టి} \textsuperscript{(49)} — దానిపై లావు \VUgras{బిందె} \textsuperscript{(50)}
పెట్టి ఉంది — తన చిన్న బీరువాతో బహుశా చాలా సౌకర్యంగా ఉంటుంది; ఆ బీరువాలో \VUgras{వెనిగర్ పీపా}
\textsuperscript{(51)} — దాని \VUgras{కుళాయితో} \textsuperscript{(52)} —,
\VUgras{చెత్త పెట్టె} \textsuperscript{(53)} మరియు \VUgras{ఎంగిలి పళ్ళేలు}
\textsuperscript{(54)} ఉంచుతారు.

%%K t06-tit-9 sub
\VUsaut{4.0mm}
\VUpk{10.2pt}{40}{వంటగదిలో ఉంచిన ఇతర వస్తువులు.}
\VUsaut{3.0mm}

%%K t06-c4-11-1 p
11. --- \VUgras{పెద్ద తొట్టి} \textsuperscript{(55)} — దానిలో \VUgras{కూరగాయలు}
\textsuperscript{(58)} కడుగుతారు — మరియు బిడ్డ ఇనుప \VUgras{గంగాళం} \textsuperscript{(56)} —
అది \VUgras{పలకల నేలపై} \textsuperscript{(57)} పెట్టి ఉంది — వాటి చోటు కూడా బహుశా ఆ బీరువాలోనే
ఉంటుంది.

%%K t06-c4-12-1 p
12. --- వంటమనిషికి పొయ్యిల కొరత ఖచ్చితంగా లేదు, ఎందుకంటే ఇదిగో ఇంకా, బల్ల మూలలో, ఒక
\VUgras{గ్యాస్ పొయ్యి} \textsuperscript{(59)}, దానిపై ఒక చిన్న గిన్నెలో నీళ్ళు
వేడెక్కుతున్నాయి. దానినుంచి వచ్చే \VUgras{ఆవిరి} \textsuperscript{(60)} అవి బహుశా
మరుగుతున్నాయని మనకు చెబుతోంది. ఈ నీళ్ళు బహుశా కాఫీకి వాడతారు, ఎందుకంటే ఇదిగో బల్ల ఇంకో చివర
\VUgras{కాఫీ కూజా} \textsuperscript{(64)} మరియు \VUgras{కాఫీ మిల్లు} \textsuperscript{(63)}.

%%K t06-c4-13-1 p
13. --- పొయ్యి \VUgras{గ్యాస్ మంటకు} \textsuperscript{(62)} కలిపి ఉంది, పొడవాటి
\VUgras{రబ్బరు గొట్టంతో} \textsuperscript{(61)}.

%%K t06-c4-14-1 p
14. --- వంటమనిషికి సాయం చేసేందుకు వచ్చే \VUgras{రోజువారీ పనిమనిషి} \textsuperscript{(66)}
భోజనానికి కూరగాయలు సిద్ధం చేస్తోంది. అవి కడిగి శుభ్రం చేశాక వాటిని \VUgras{పెద్ద పళ్ళెంలో}
\textsuperscript{(65)} పెట్టి వండే సమయం కోసం వేచి ఉంటుంది.

%%K t06-tit-10 sub
\VUsaut{4.0mm}
{\centering\fontsize{10.2pt}{10.2pt}\selectfont\textls[40]{\scshape స్నానపు గది.}
\textit{(స్నానగృహం.)}\par}
\VUsaut{3.0mm}

%%K t06-c5-01-1 p
1. --- ఇంట్లోని ముఖ్యమైన గదులు తెలుసుకోవాలంటే ఇక ఒక్క అమర్యాద మిగిలింది : స్నానపు గది అమరిక
చూడటం. దానికి ఎక్కువ సమయం పట్టదు, ఎందుకంటే అది పెద్దది కాదు; \VUgras{స్నానపు తొట్టికి}
\textsuperscript{(1)} మంచి \VUgras{నీళ్ళు కాచే యంత్రం} \textsuperscript{(2)} ఉందని,
\VUgras{సబ్బు పెట్టె} \textsuperscript{(3)} స్నానం చేసేవాడి చేతికి అందేలా ఉందని,
\VUgras{తువ్వాళ్ళ దండెం} \textsuperscript{(5)} \VUgras{తువ్వాళ్ళను} \textsuperscript{(4)}
సులువుగా ఆరబెడుతుందని మనకు వెంటనే తెలుస్తుంది.

%%K t06-c5-02-1 p
2. --- ఈ గది గోడలకు \VUgras{పింగాణీ పలకలు} \textsuperscript{(6)} వేసి ఉన్నాయి, వాటివల్లే
\VUgras{జల్లు స్నానం} \textsuperscript{(7)} అమర్చడం సాధ్యమైంది — వాడేటప్పుడు చిమ్మే నీళ్ళు
గోడలు పాడుచేస్తాయనే భయం లేకుండా. కాళ్ళ స్నానానికి వాడే జింక్ \VUgras{పళ్ళెం}
\textsuperscript{(8)} ఈ సామానును పూర్తి చేస్తుంది.

\VUsaut{6.0mm}
\VUfilet{22mm}
